\section{O corpo dos números reais}
A estrutura dos números reais é caracterizada por consistir em um conjunto, denotado por $\mathbb R$, munido de operações binárias $+$, $\cdot$, elementos destacados $0$ e $1$, e uma relação $<$ que satisfaz certas propriedades.

Assumiremos, sem prova, que tal estrutura existe.
Abaixo, enunciaremos todas as propriedades que caracterizam esta estrutura.
A seguir, estudaremos-a em detalhes.

\begin{definition}
    A estrutura dos números reais consiste em um conjunto $\mathbb R$ munido de operações binárias $+: \mathbb R\times \mathbb R\to \mathbb R$ e $\cdot: \mathbb R\times \mathbb R\to \mathbb R$, elementos destacados $0$ e $1$, e uma relação $<\subseteq \mathbb R\times \mathbb R$ tal que, para quaisquer $x, y, z\in \mathbb R$:
    \begin{enumerate}
        \labeleditem{(A1)} $x+(y+z)=(x+y)+z$ (propriedade associativa);\label{defItem:associativaSoma} \label{defItem:primeiraSoma}
        \labeleditem{(A2)} $x+y=y+x$ (propriedade comutativa); \label{defItem:comutativaSoma}
        \labeleditem{(A3)} $x+0=x$ (elemento neutro para adição); \label{defItem:neutroSoma}
        \labeleditem{(A4)} existe $u\in \mathbb R$ tal que $x+u=0$ (existência de elemento inverso para adição); \label{defItem:inversoSoma} \label{defItem:ultimoSoma}
        \labeleditem{(M1)} $x\cdot (y\cdot z)=(x\cdot y)\cdot z$ (propriedade associativa); \label{defItem:associativaMultiplicacao} \label{defItem:primeiraMultiplicacao}
        \labeleditem{(M2)} $x\cdot 1=x$ e $1\cdot x=x$ (elemento neutro para multiplicação); \label{defItem:neutroMultiplicacao}
        \labeleditem{(M3)} $x(y+z)=x\cdot y+x\cdot z$ (propriedade distributiva); \label{defItem:distributiva}
        \labeleditem{(M4)} Se $x\neq 0$, existe $v\in \mathbb R$ tal que $x\cdot v=1$ e $v\cdot x=1$ (existência de elemento inverso para multiplicação); \label{defItem:inversoMultiplicacao}
        \labeleditem{(M5)} $0\neq 1$ (não trivialidade); \label{defItem:naoTrivialidade} \label{defItem:ultimoMultiplicacao}
        \labeleditem{(O1)} Se $x<y$ e $y<z$, então $x<z$ (propriedade transitiva); \label{defItem:primeiraOrdem} \label{defItem:transitiva}
        \labeleditem{(O2)} $x \not< x$ (propriedade irreflexiva); \label{defItem:segundaOrdem} \label{defItem:irreflexiva}
        \labeleditem{(O3)} $x<y$ ou $x=y$ ou $y<x$ (propriedade de totalidade); \label{defItem:terceiraOrdem} \label{defItem:totalidade}
        \labeleditem{(OA)} Se $x<y$, então $x+z<y+z$ (compatibilidade da adição com a ordem); \label{defItem:compatibilidadeSoma}
        \labeleditem{(OM)} Se $x<y$ e $0<z$, então $x\cdot z<y\cdot z$ (compatibilidade da multiplicação com a ordem); \label{defItem:compatibilidadeMultiplicacao}
        \labeleditem{(S)} Todo subconjunto de $\mathbb R$ não vazio e limitado superiormente tem um supremo (propriedade de completude). \label{defItem:supremo}
    \end{enumerate}

    Os elementos de $\mathbb R$ são chamados de números reais.
\end{definition}

A Propriedade~\ref{defItem:supremo}, também conhecida como a propriedade do supremo, certamente requer explicações adicionais.
Deixá-la-emos por último.
Solicitamos que o leitor não se preocupe com ela até que tal explicação seja concedida.

Iniciaremos discutindo as propriedades do grupo \ref{defItem:primeiraSoma}-\ref{defItem:ultimoSoma}.
Tais propriedades regem a adição de números reais isoladamente.
Ela nos permite definir a operação de oposto.

\begin{proposition}
    Seja $x$ um número real.
    Então existe um único número real $u$ tal que $x+u=0$.
\end{proposition}
\begin{proof}
    Tome um número real $x$.
    Pela propriedade \ref{defItem:inversoSoma}, existe um número real $u$ tal que $x+u=0$.
    Vejamos que só existe um: ou seja, que se $u$ e $u'$ são números reais tais que $x+u=0$ e $x+u'=0$, então $u=u'$.

    De fato, se $u$ e $u'$ são números reais tais que $x+u=0$ e $x+u'=0$, então
\begin{align*}
    u &= u + 0 \tag{0 é neutro da soma} \\
    &= u + (x + u) \tag{pela escolha de $u$} \\
    &= (u + x) + u' \tag{propriedade associativa} \\
    &= 0 + u' \tag{pela escolha de $u'$} \\
    &= u' \tag{0 é neutro da soma}
\end{align*}
Assim, $u=u'$.
\end{proof}

Dado um número real $x$, existe então um único número real $u$ tal que $x+u=0$.
Tal número $u$ é completamente determinado por essa propriedade, e depende exclusivamente de $x$.
Nada mais justo do que dar a ele um nome conveniente.

\begin{definition}
    Seja $x$ um número real.
    O \emph{oposto}\index{oposto} de $x$, denotado por $-x$, é o único número real tal que $x+(-x)=0$.
\end{definition}

Com isso, podemos provar algumas das chamadas \emph{regras de sinais}.
Primeiro, uma convenção.

\begin{definition}
    Sejam $x$, $y$ números reais.
    Define-se a operação de subtração de números reais $-: \mathbb R\times \mathbb R\to \mathbb R$ por $x-y = x+(-y)$.
\end{definition}

Percebamos algo importante: há uma ambiguidade na notação de $-$.
Acabamos de definir a notação $x-y$ como uma operação binária: $-(x, y)=x-y=x+(-y)$.
Tal operação requer dois elementos.
Ao mesmo tempo, a notação $-$ remete a um operador unário: dado $x$, $-x$ é o oposto de $x$.

Utilizando as convenções aritméticas usuais, acaba não persistindo ambiguidade: se $-$ é escrito sem nada à esquerda que esteja claramente operando, então $-$ denota o oposto.
Caso contrário, $-$ denota a operação de subtração.
Parenteses são usados para evitar ambiguidades.

\begin{proposition}
    Sejam $x, y$ números reais.
    Então:
    \begin{enumerate}[label=(\alph*)]
        \item $-0=0$;
        \item $-(x+y)=(-x)+(-y)$;
        \item $-(-x)=x$.
    \end{enumerate}
\end{proposition}
\begin{proof}
    Para a propriedade $(a)$, temos que $0+0=0$. Como $-0$ é o único número real tal que $0+(-0)=0$, então $-0=0$.

    Para a propriedade $(b)$, temos que
    \begin{align*}
    (x+y)+((-x)+(-y)) &= [(x+y)+(-x)]+(-y) \tag{propriedade associativa} \\
    &= [(y+x)+(-x)]+(-y) \tag{propriedade comutativa} \\
    &= [y+(x+(-x))]+(-y) \tag{propriedade associativa} \\
    &= (y+0)+(-y) \tag{pela definição de oposto} \\
    &= y+(-y) \tag{0 é neutro da soma} \\
    &= 0. \tag{pela definição de oposto}
    \end{align*}
    Como $-(x+y)$ é o único número real tal que $(x+y)+(-(x+y))=0$, então $-(x+y)=(-x)+(-y)$.

    Para a propriedade $(c)$, temos que $(-x)+x=0$.
    Como $-(-x)$ é o único número real tal que $(-x)+(-(-x))=0$, então $-(-x)=x$.
\end{proof}

Notemos que a propriedade associativa nos diz que se $x, y, z$ são números reais, então $x+(y+z)=(x+y)+z$.
Assim, não há ambiguidade em escrever $x+y+z$ sem parênteses.
No futuro, escreveremos expressões como essa sem ressalvas.

Outras consequências dessas propriedades são as seguintes.
\begin{proposition}
    Seja $x$ um número real.
    Se $-x=0$, então $x=0$.
    Em particular, $-1\neq 0$.
\end{proposition}
\begin{proof}
    Se $-x=0$, então $x=(-(-x))=-0=0$, provando a primeira afirmação.
    Para a segunda, aplique a primeira com $x=1$.
    Como $1\neq 0$, segue que $-1\neq 0$ (ou teríamos $-1=0$, nos dando $1=0$).
\end{proof}

Agora vamos discutir um pouco acerca das propriedades do grupo \ref{defItem:primeiraMultiplicacao}-\ref{defItem:ultimoMultiplicacao}.
Tais propriedades regem a multiplicação de números reais e sua relação com a soma.

Começaremos mostrando que o produto por $0$ é sempre $0$.

\begin{proposition}
Seja $x$ um número real.
Então $x\cdot 0=0$.
\end{proposition}
\begin{proof}
    Temos que:
    \begin{align*}
    x\cdot 0 &= x\cdot (0+0) \tag{0 é neutro da soma} \\
    &= x\cdot 0 + x\cdot 0 \tag{propriedade distributiva} \\
    \end{align*}

    Assim, $x\cdot 0 = x\cdot 0 + x\cdot 0$.
    Temos, somando a ambos os lados $-(x\cdot 0)$, temos:
    \begin{align*}
    0 &= -(x\cdot 0) + (x\cdot 0 + x\cdot 0) \tag{pela propriedade do oposto $-(x\cdot 0)$} \\
    &= (-(x\cdot 0) + x\cdot 0) + x\cdot 0 \tag{propriedade associativa} \\
    &= 0 + x\cdot 0 \tag{pela definição de oposto} \\
    &= x\cdot 0. \tag{0 é neutro da soma}
    \end{align*}
    Portanto, $x\cdot 0=0$.
\end{proof}

Podemos também demonstrar as regras de sinais que dizem respeito à multiplicação.

\begin{proposition}
    Sejam $x, y$ números reais.
    Então:
    \begin{enumerate}[label=(\alph*)]
        \item $(-x)\cdot y = -(x\cdot y)=x\cdot (-y)$;
        \item $(-x)(-y)=x\cdot y$.
        \item $-x = (-1)\cdot x$.
    \end{enumerate}
\end{proposition}
\begin{proof}
    Para a primeira propriedade, notemos que $x\cdot y + (-(x\cdot y))=x\cdot[y+(-y)]=x\cdot 0=0$.
    Como $-(x\cdot y)$ é o único número real tal que $(x\cdot y) + (-(x\cdot y))=0$, então $(-x)\cdot y = -(x\cdot y)$.

    Para a segunda igualdade do primeiro item, note que, usando a primeira igualdade, que já provamos, temos que $x\cdot(-y)=(-y)\cdot x=-(y\cdot x)=-(x\cdot y)$.

    Para o segundo item, iteramos a primeira propriedade:

    \[
    (-x)(-y) = -(x\cdot(-y)) = -(-(x\cdot y)) = x\cdot y.
    \]

    Para o último item, do primeiro, temos que $(-1)\cdot x = 1\cdot (-x) = -x$.
\end{proof}

Outra propriedade fundamental dos números reais é a seguinte.
\begin{definition}
    Suponha que $x$ e $y$ são números reais tais que $x\cdot y=0$.
    Então $x=0$ ou $y=0$.
\end{definition}
\begin{proof}
    Suponha que $x\cdot y=0$.
    Se $x=0$, não há nada para provar.
    Se $x\neq 0$, então existe um número real $v$ tal que $v\cdot x=1$.
    Assim, temos que
    \[
    y=1\cdot y=(v\cdot x)\cdot y=v\cdot (x\cdot y)=v\cdot 0=0.
    \]
\end{proof}
Abaixo, temos a proposição que nos permite definir frações da forma $\frac{a}{b}$.

\begin{proposition}
    Sejam $a, b$ números reais com $b\neq 0$.
    Existe um único número real $r$ tal que $b\cdot r=a$.
\end{proposition}

\begin{proof}
    Para a existência, pela propriedade \ref{defItem:inversoMultiplicacao}, existe um número real $v$ tal que $b\cdot v=1$.
    Tomando $r=v\cdot a$, temos que:
    \begin{align*}
    b\cdot r &= b\cdot (v\cdot a) \tag{pela definição de $r$} \\
    &= (b\cdot v)\cdot a \tag{propriedade associativa} \\
    &= 1\cdot a \tag{pela escolha de $v$} \\
    &= a. \tag{elemento neutro da multiplicação}
    \end{align*}

    Para a unicidade, sejam $r$ e $r'$ números reais tais que $b\cdot r=a$ e $b\cdot r'=a$.
    Tome $v$ um número real tal que $b\cdot v=1$.
    
    Temos que $b\cdot r = b\cdot r'$.
    Dessa forma, segue que
    \[
        v \cdot (b\cdot r) = v \cdot (b\cdot r').
    \]

    Pela propriedade associativa, temos
    \[
        (v \cdot b)\cdot r = (v \cdot b)\cdot r'.
    \]
    
    Pela escolha de $v$, temos
    \[
        1\cdot r = 1\cdot r'.
    \]
    
    Como $1$ é elemento neutro, temos
    \[
        r = r',
    \]
    como queríamos mostrar.
\end{proof}

\begin{definition}
    Sejam $a, b$ números reais com $b\neq 0$.
    A \emph{fração}\index{fração} $\dfrac{a}{b}$ é o único número real que satisfaz

    \[
        b\cdot \frac{a}{b} = a.
    \]
\end{definition}

Notemos que, não é possível definir $\frac{a}{0}$.
Se para algum $a \in \mathbb R\setminus\{0\}$ houvesse um número real $r$ tal que $0\cdot r=a$, então teríamos que $0=0\cdot r=a$.
Mesmo que $a=0$, não é uma boa prática definir $\frac{0}{0}$.
Apesar de ser verdade que existe $u$ tal que $0\cdot u=0$ (por exemplo, $u=0$), tal $u$ não é único: qualquer número real satisfaz $0\cdot u=0$ (por exemplo, $0$ e $1$ são distintos e satisfazem essa equação).

Agora, demonstremos propriedades referentes às frações.
A partir deste ponto, omitiremos frequentemente o símbolo de multiplicação $\cdot$, escrevendo apenas $ab$ ao invés de $a\cdot b$, a fim de limpar a notação.

\begin{proposition}
    Sejam $a, b, c, d$ números reais, com $b\neq 0$ e $d\neq 0$ (de modo que, consequentemente, $bd\neq 0$).
    Então:
    \begin{enumerate}[label=(\alph*)]
        \item $\dfrac{a}{b}\cdot \dfrac{c}{d} = \dfrac{ac}{bd}$;
        \item $\dfrac{d}{d}=1$;
        \item $\dfrac{c}{1} = c$;
        \item $\dfrac{a}{b} = \dfrac{ad}{bd}$;
        \item $\dfrac{a}{b} + \dfrac{c}{b} = \dfrac{a+c}{b}$;
        \item $\dfrac{a}{b} + \dfrac{c}{d} = \dfrac{ad + bc}{bd}$;
        \item $\dfrac{ac}{b} = a\cdot \dfrac{c}{b}$;
        \item $- \dfrac{a}{b} = \dfrac{-a}{b} = \dfrac{a}{-b}$;
        \item $\dfrac{a}{c} - \dfrac{b}{c} = \dfrac{a-b}{c}$;
        \item $\dfrac{a}{b} - \dfrac{c}{d} = \dfrac{ad - bc}{bd}$;
        \item Se $a\neq 0$, então $\dfrac{a}{b}\neq 0$.
        \item Se $c\neq 0$, então $\dfrac{\dfrac{a}{b}}{\dfrac{c}{d}} = \dfrac{ad}{bc}$.
    \end{enumerate}
\end{proposition}
\begin{proof}\hfill
    (a) Basta ver que $\left(\dfrac{a}{b}\cdot \dfrac{c}{d}\right)\cdot bd = ac$.
    De fato, pelas propriedades comutativa e associativa da multiplicação, temos que
    \[
    \left(\dfrac{a}{b}\cdot \dfrac{c}{d}\right)\cdot bd = \left(\dfrac{a}{b} \cdot a\right) \cdot \left(\dfrac{c}{d}\cdot d\right) = ac.
    \]

    (b) Imediato, uma vez que $d\cdot 1=  d$.

    (c) Imediato, uma vez que $1\cdot c = c$.

    (d) Dos itens anteriores, temos que:
    \[
    \dfrac{ad}{bd} = \dfrac{a}{b}\cdot \dfrac{d}{d} = \dfrac{a}{b}\cdot 1 = \dfrac{a}{b}.
    \]

    (e) Basta ver que $\left(\dfrac{a}{b} + \dfrac{c}{b}\right)\cdot b = a+c$.
    De fato, pela propriedade distributiva, temos que
    \[
    \left(\dfrac{a}{b} + \dfrac{c}{b}\right)\cdot b = \dfrac{a}{b}\cdot b + \dfrac{c}{b}\cdot b = a+c.
    \]

    (f) Das propriedades anteriores, temos que
    \[
        \dfrac{a}{c} + \dfrac{b}{d} = \dfrac{ad}{cd} + \dfrac{bc}{cd} = \dfrac{ad + bc}{cd}.
    \]

    (g) Das propriedades anteriores, temos que
    \[
        \dfrac{ac}{b} = \dfrac{a\cdot c}{1\cdot b} = \dfrac{a}{1}\cdot \dfrac{c}{b} = a\cdot \dfrac{c}{b}.
    \]
    (h) Das propriedades anteriores, temos que
    \[
        -\dfrac{a}{b} = (-1)\cdot \dfrac{a}{b} = \dfrac{(-1)a}{b} = \dfrac{-a}{b}
    \]
    Note que a fração $\dfrac{a}{-b}$ faz sentido, uma vez que $-b=0$ implica $b=0$, o que é uma contradição.
    
    Assim, $\dfrac{a}{-b}$ é bem definida, e temos que
    \[
        \dfrac{a}{-b} = \dfrac{a\cdot (-1)}{(-b)(-1)} = \dfrac{-a}{b}.
    \]

    (i) Pelas propriedades anteriores, temos que
    \[
        \dfrac{a}{c} - \dfrac{b}{c} = \dfrac{a}{c} + (-1)\cdot \dfrac{b}{c} = \dfrac{a}{c} + \dfrac{-b}{c} = \dfrac{a+(-b)}{c}=\dfrac{a-b}{c}.
    \]

    (j) Pelas propriedades anteriores, temos que
    \[
        \dfrac{a}{b} - \dfrac{c}{d} = \dfrac{ad}{bd} - \dfrac{bc}{bd} = \dfrac{ad - bc}{bd}.
    \]

    (k) Suponha que $\frac{a}{b}=0$.
    Então $a=b\frac{a}{b}=b\cdot 0=0$
    Dessa forma, se $a\neq 0$, então $\frac{a}{b}\neq 0$ (do contrário, teríamos $\frac{a}{b}=0$, o que implica $a=0$).

    (l) Suponha que $c\neq 0$.
    Então $\dfrac{c}{d}\neq 0$ e $bc\neq 0$.

    Basta ver que $\left(\dfrac{ad}{bc}\right)\cdot \dfrac{c}{d} = \dfrac{a}{b}$.
    De fato, pelas propriedades comutativa e associativa da multiplicação, temos que
    \[
        \left(\dfrac{ad}{bc}\right)\cdot \dfrac{c}{d} = \left(\dfrac{a}{b}\cdot \dfrac{d}{c}\right)\cdot \dfrac{c}{d} = \dfrac{a}{b}\cdot \left(\dfrac{d}{c}\cdot \dfrac{c}{d}\right) = \dfrac{a}{b}\cdot 1 = \dfrac{a}{b}.
    \]
\end{proof}

Agora procederemos para o estudo das propriedades sobre a ordem $<$.
A Propriedade~\ref{defItem:totalidade} nos diz que quaisquer dois números reais $x$ e $y$ são comparáveis: ou eles são o mesmo, ou um deles é o maior.
Porém, a Propriedade~\ref{defItem:totalidade} não nos diz que apenas uma dessas possibilidades ocorre.
Tal fato é, porém, uma consequência das outras propriedades.

\begin{proposition}
    Sejam $x$ e $y$ números reais.
    Os seguintes itens NÃO podem ocorrer:
    \begin{enumerate}[label=(\alph*)]
        \item $x<y$ e $y<x$;
        \item $x<y$ e $x=y$;
        \item $y<x$ e $x=y$.
    \end{enumerate}
\end{proposition}
\begin{proof}
    Suponha que $x<y$ e $y<x$.
    Pela propriedade transitiva, $x<y$ e $y<x$ implicam $x<x$, o que viola a propriedade irreflexiva.

    Suponha que $x<y$ e $x=y$.
    Substituindo $y$ por $x$ na primeira desigualdade, temos que $x<x$, o que viola a propriedade irreflexiva.

    Por fim, suponha que $y<x$ e $x=y$.
    Substituindo $x$ por $y$ na primeira desigualdade, temos que $y<y$, o que viola a propriedade irreflexiva.
\end{proof}

Símbolos muito usados em Matemática são as relações de ordem não estritas: $\leq$.

\begin{definition}
    Sejam $x$ e $y$ números reais.
    Dizemos que:
    \begin{itemize}
        \item $x$ é \emph{menor ou igual a} $y$, denotado por $x\leq y$, se $x<y$ ou $x=y$;
        \item $x$ é \emph{maior} que $y$, denotado por $x>y$, se $y<x$;
        \item $x$ é \emph{maior ou igual a} $y$, denotado por $x\geq y$, se $y\leq x$.
    \end{itemize}
\end{definition}

\begin{proposition}
    Sejam $x$, $y$ e $z$ números reais.
    \begin{enumerate}[label=(\alph*)]
        \item Se $x\leq y$ e $y\leq z$, então $x\leq z$;
        \item Se $x\leq y$ e $y\leq x$, então $x=y$;
        \item $x\leq x$.
        \item Se $x\leq y$, então $x+z \leq y+z$.
        \item Se $x\leq y$ e $0\leq z$, então $xz \leq yz$.
    \end{enumerate}
\end{proposition}
\begin{proof}\hfill
    (a) Suponha que $x\leq y$ e $y\leq z$.
    
    Se $x=y$, temos que $y\leq z$ substituindo $y$ por $x$ na segunda desigualdade.

    Se $y=z$, temos que $x\leq y$ substituindo $y$ por $z$ na primeira desigualdade.

    Caso contrário, temos que $x<y$ e $y<z$, o que implica $x<z$ pela propriedade transitiva.

    (b) Suponha que $x\leq y$ e $y\leq x$.

    Suponha por absurdo que $x\neq y$.
    Segue da definição de $\leq$ que $x<y$ e $y<x$, o que é uma contradição.

    (c) Temos que $x=x$.
    Logo, pela definição de $\leq$, $x\leq x$.

    (d) Suponha que $x\leq y$.
    Se $x=y$, então $y+z=x+z$ substituindo-se $y$ por $x$.
    Se $x<y$, então $x+z<y+z$ pela propriedade de compatibilidade da adição com a ordem.

    (e) Suponha que $x\leq y$ e $0\leq z$.
    Se $x=y$, segue que $xz=yz$, e, portanto, $xz\leq yz$.
    A partir daqui, suponha então que $x<y$.

    Se $0=z$, temos que $xz=x\cdot 0 = 0$ e $yz=y\cdot 0=0$, de modo que $xz\leq yz$.

    Por fim, resta o caso em que $0<z$ e $x<y$.
    Pela propriedade de compatibilidade da multiplicação com a ordem, $xz<yz$, o que implica $xz\leq yz$.
\end{proof}

Isso será provado abaixo.
É importante notar que o símbolo $-x$ não indica que $x$ é negativo.
A operação $x\mapsto -x$ é apenas uma operação algébrica que mapeia $x$ em seu oposto aditivo.
Note que para $x=-1$, temos que $-x=-(-1)=1$, que é positivo.
Em verdade, ainda não provamos que $1>0$.
Isso será feito abaixo.

Denotamos, por $x^2$, o produto $x\cdot x$.

\begin{proposition}
    Para todo número real $x$, temos que $x^2\geq 0$.
    Além disso, $x^2=0$ se e somente se $x=0$.
\end{proposition}
\begin{proof}
    Fixe um número real $x$.
    Vamos dividir a prova em casos.
    Sabemos, pela totalidade, que $x<0$ ou $x=0$ ou $0<x$.

    Caso $0<x$, temos que $0<x$ e $0<x$, de modo que, pela propriedade de compatibilidade da multiplicação com a ordem, $0x <xx=x^2$.
    Assim, caso $0<x$, temos que $0<x^2$.

    Caso $x<0$, temos que, pela compatibilidade da adição com a ordem, $x+(-x)<0+(-x)$, ou seja, $0<-x$.
    Assim, pelo caso anterior, $0<(-x)^2=(-x)(-x)=x^2$.

    Por fim, caso $x=0$, temos que $x^2=0^2=0\cdot 0=0$.

    Notemos que o único destes casos em que $x^2=0$ é o caso em que $x=0$.
\end{proof}

Agora vamos provar nossas primeiras desigualdades estritas concretas.
\begin{proposition}
    Nos números reais, $-1<0<1$.
\end{proposition}
\begin{proof}
    Note que $1\cdot 1=1$, logo, $1=1^2\geq 0$.
    Como $1\neq 0$, segue que $0<1$.
    Somando $(-1)$ a ambos os lados da desigualdade $0<1$, temos que $-1<0$.

    Assim, $-1<0<1$.
\end{proof}